% Plan de habilitación · Filosofía 10.° · Periodo IV
% Para el estudiante que no alcanzó uno o más desempeños del periodo (D1–D5 de malla-P4.tex).
% Funciona en dos pasos: (1) trabajo de habilitación en casa sobre los desempeños marcados;
% (2) evaluación de habilitación (skill evaluaciones-latex) sobre esos mismos desempeños.
% Resultado: aprobado o no aprobado (criterios definidos por el docente; decisión del docente, nov. 2026).
% Formato oficial: el mismo encabezado institucional con el nombre del formato cambiado.
% El docente marca a mano las casillas de los desempeños que el estudiante debe recuperar.
% Compilar: bash curriculo/plantilla/compilar.sh curriculo/filosofia/10/habilitacion-P4.tex
\documentclass[habilitacion]{planeacion}
\materia{filosofia}
\grado{10.°} \periodo{IV} \anio{2027}
\carpeta{}
\begin{document}
\fichahabilitacion

\destacado{CÓMO\newline FUNCIONA}{Primero entregas un \textbf{trabajo} sobre los aprendizajes
  marcados; después presentas una \textbf{evaluación} corta en la que demuestras que ya los
  alcanzaste. El trabajo es tu preparación: si lo haces bien, la evaluación no tendrá sorpresas.}

\seccion{1}{Lo que debes recuperar}
El docente marca (\casilla) los aprendizajes que debes recuperar. Haz \textbf{solo} las actividades
de los marcados, a mano, en hojas aparte o en tu cuaderno, con tu nombre en cada hoja.

\begin{recuperacion}
\recuperar{D1}{Sabes cómo llegaron las ideas de Platón y Aristóteles a la Edad Media.}
  {Línea de tiempo con Platón, Aristóteles, Agustín, la caída de Roma, las primeras universidades,
  Tomás de Aquino y la imprenta. Debajo, en media página, el «viaje» de los libros de Aristóteles:
  quiénes los conservaron y tradujeron, y dónde. Menciona al menos dos cosas que hicimos en clase
  (el cuadro de Rafael, el reto de traducción).}
  {apuntes firmados de la semana 1 y el mapa del viaje de los textos.}
\recuperar{D2}{Interpretas a San Agustín y su idea de la relación entre fe y razón.}
  {Explica con tus palabras «comprende para creer, cree para comprender». Luego, un párrafo: ¿estás
  de acuerdo con Agustín? Da una razón y un ejemplo de tu propia vida.}
  {fragmentos de Agustín (carpeta del periodo) y tu carta a Agustín.}
\recuperar{D3}{Conoces el método de la \emph{quaestio} y sabes usarlo.}
  {Una \emph{quaestio} completa sobre una pregunta distinta a la de tu taller: pregunta, dos
  objeciones, «en sentido contrario», tu respuesta y la respuesta a cada objeción.}
  {el modelo del Taller 1 y el esquema de la disputa.}
\recuperar{D4}{Evalúas con evidencias si la Edad Media fue una «época oscura» y explicas qué cambió
  con la imprenta y el humanismo.}
  {Un asistente de IA respondió: \textit{«La Edad Media fue una época oscura: la Iglesia prohibió
  leer a los filósofos griegos, no existían escuelas ni universidades y todo cambió de un día para
  otro cuando Gutenberg inventó la imprenta.»} Señala \textbf{tres errores} y corrígelos con
  evidencias de clase. Termina con tu conclusión y un párrafo sobre qué cambió con la imprenta y
  con Erasmo.}
  {veredicto del juicio, cálculo de la carrera copista contra imprenta y cuadro comparativo.}
\recuperar{D5}{Interpretas la crítica social de Tomás Moro y propones una utopía con sus riesgos.}
  {Explica qué quiere decir Moro con las «ovejas que se comen a los hombres». Luego propone una ley
  para un problema real del municipio o del colegio, con su razón y un riesgo.}
  {fragmentos de la \emph{Utopía} y el Taller 2.}
\end{recuperacion}

\seccion{2}{Entrega, evaluación y resultado}
\begin{itemize}\setlength{\itemsep}{1pt}
  \item \textbf{Trabajo:} a mano, completo y en la fecha indicada arriba. Sin trabajo no se presenta
    la evaluación.
  \item \textbf{Sobre la IA:} puedes consultar, incluso con inteligencia artificial, pero el trabajo
    va a mano, con tus palabras y con referencias a tus apuntes. Si usaste IA, escribe al final para
    qué la usaste. La evaluación es presencial: ahí se ve lo que de verdad aprendiste.
  \item \textbf{Evaluación de habilitación:} escrita y corta, solo sobre los aprendizajes marcados,
    con el mismo tipo de preguntas de la evaluación de periodo (textos, imágenes y argumentos).
  \item \textbf{Resultado:} la habilitación se \textbf{aprueba} o \textbf{no se aprueba}; no tiene
    nota intermedia.
\end{itemize}

\seccion{3}{Para aprobar}
Apruebas si cumples las tres condiciones (el docente las marca):

\medskip
\begin{tabular}{@{}l p{0.86\linewidth}@{}}
  \casilla & Entregaste el trabajo \textbf{completo y a tiempo}: todas las actividades marcadas,
             hechas a mano y explicadas con tus palabras (no copiadas).\\[3pt]
  \casilla & En el trabajo \textbf{das razones y ejemplos propios}, no solo opiniones.\\[3pt]
  \casilla & En la evaluación alcanzas al menos el \textbf{nivel básico} (60\,\% de los puntos) en
             \textbf{cada} aprendizaje marcado.\\
\end{tabular}

\medskip
\noindent\textbf{Resultado} (lo marca el docente):\quad \casilla~Aprobado\qquad \casilla~No aprobado

\recibido
\end{document}
